\documentclass[10pt,a4paper]{amsart}
\usepackage{amsmath,amssymb,amsthm,graphicx}
% Retained minimal preamble for the Stacks Project chapter sites-cohomology.tex (Cohomology on Sites).
% Verbatim excerpt of the chapter preamble: Xy-pic setup (native preamble lines 14-18) and the native
% theorem declarations with the mathematical macros (native preamble lines 154-217). Copied mechanically;
% nothing added, removed or redefined. The wrapper class and the amsmath/amssymb/amsthm/graphicx base
% packages are supplied by the assembler, so the class/externaldocument part of the native preamble is omitted.
% For commutative diagrams we use Xy-pic
\usepackage[all]{xy}

% We use 2cell for 2-commutative diagrams.
\xyoption{2cell}
% Theorem environments.
%
\theoremstyle{plain}
\newtheorem{theorem}[subsection]{Theorem}
\newtheorem{proposition}[subsection]{Proposition}
\newtheorem{lemma}[subsection]{Lemma}

\theoremstyle{definition}
\newtheorem{definition}[subsection]{Definition}
\newtheorem{example}[subsection]{Example}
\newtheorem{exercise}[subsection]{Exercise}
\newtheorem{situation}[subsection]{Situation}

\theoremstyle{remark}
\newtheorem{remark}[subsection]{Remark}
\newtheorem{remarks}[subsection]{Remarks}

\numberwithin{equation}{subsection}

% Macros
%
\def\lim{\mathop{\mathrm{lim}}\nolimits}
\def\colim{\mathop{\mathrm{colim}}\nolimits}
\def\Spec{\mathop{\mathrm{Spec}}}
\def\Hom{\mathop{\mathrm{Hom}}\nolimits}
\def\Ext{\mathop{\mathrm{Ext}}\nolimits}
\def\SheafHom{\mathop{\mathcal{H}\!\mathit{om}}\nolimits}
\def\SheafExt{\mathop{\mathcal{E}\!\mathit{xt}}\nolimits}
\def\Sch{\mathit{Sch}}
\def\Mor{\mathop{\mathrm{Mor}}\nolimits}
\def\Ob{\mathop{\mathrm{Ob}}\nolimits}
\def\Sh{\mathop{\mathit{Sh}}\nolimits}
\def\NL{\mathop{N\!L}\nolimits}
\def\CH{\mathop{\mathrm{CH}}\nolimits}
\def\proetale{{pro\text{-}\acute{e}tale}}
\def\etale{{\acute{e}tale}}
\def\QCoh{\mathit{QCoh}}
\def\Ker{\mathop{\mathrm{Ker}}}
\def\Im{\mathop{\mathrm{Im}}}
\def\Coker{\mathop{\mathrm{Coker}}}
\def\Coim{\mathop{\mathrm{Coim}}}

% Boxtimes
%
\DeclareMathSymbol{\boxtimes}{\mathbin}{AMSa}{"02}

%
% Macros for moduli stacks/spaces
%
\def\QCohstack{\mathcal{QC}\!\mathit{oh}}
\def\Cohstack{\mathcal{C}\!\mathit{oh}}
\def\Spacesstack{\mathcal{S}\!\mathit{paces}}
\def\Quotfunctor{\mathrm{Quot}}
\def\Hilbfunctor{\mathrm{Hilb}}
\def\Curvesstack{\mathcal{C}\!\mathit{urves}}
\def\Polarizedstack{\mathcal{P}\!\mathit{olarized}}
\def\Complexesstack{\mathcal{C}\!\mathit{omplexes}}
% \Pic is the operator that assigns to X its picard group, usage \Pic(X)
% \Picardstack_{X/B} denotes the Picard stack of X over B
% \Picardfunctor_{X/B} denotes the Picard functor of X over B
\def\Pic{\mathop{\mathrm{Pic}}\nolimits}
\def\Picardstack{\mathcal{P}\!\mathit{ic}}
\def\Picardfunctor{\mathrm{Pic}}
\def\Deformationcategory{\mathcal{D}\!\mathit{ef}}

% --- appended 2026-10-12 (Chief): the project own definitions of the reference/slogan/history
% display environments, copied mechanically from the shared Stacks preamble lines 8-12
% (\usepackage{verbatim} plus the three \newenvironment lines, which the chapter preamble keeps
% verbatim from the project preamble). Nothing added, removed or redefined: these environments
% carry no mathematical content and the project itself typesets them as comments. Without these
% three lines a counted statement whose own text uses them fails with Environment slogan undefined.
\usepackage{verbatim}
\newenvironment{reference}{\comment}{\endcomment}
\newenvironment{slogan}{\comment}{\endcomment}
\newenvironment{history}{\comment}{\endcomment}
\usepackage{cleveref}
\crefname{originalreference}{original reference}{original references}
\Crefname{originalreference}{Original reference}{Original references}
\makeatletter
\newlabel{situation-olsson-laszlo}{{0D6R}{source}{}{sourcepointer-1}{}}
\newlabel{situation-olsson-laszlo@cref}{{[originalreference][0][]0D6R}{[1][0][]source}{}{}{}}
\makeatother
\crefname{originallemma}{original lemma}{original lemmas}
\Crefname{originallemma}{Original lemma}{Original lemmas}
\makeatletter
\newlabel{lemma-equivalence-bounded}{{0D7U}{source}{}{sourcepointer-2}{}}
\newlabel{lemma-equivalence-bounded@cref}{{[originallemma][0][]0D7U}{[1][0][]source}{}{}{}}
\makeatother
\crefname{originallemma}{original lemma}{original lemmas}
\Crefname{originallemma}{Original lemma}{Original lemmas}
\makeatletter
\newlabel{lemma-olsson-laszlo-map-version-one}{{0D6U}{source}{}{sourcepointer-3}{}}
\newlabel{lemma-olsson-laszlo-map-version-one@cref}{{[originallemma][0][]0D6U}{[1][0][]source}{}{}{}}
\makeatother
\title{Equivalence of unbounded derived categories of suitable subcategories, in the first variant}
\author{The Stacks Project Authors}
\date{stacks-project:sites-cohomology.tex:local-snapshot-20261011; GNU FDL 1.2 or later}
\begin{document}
\maketitle
\noindent\textit{Editorial note.} All delivered mathematical text is the author's own, copied from the pinned Stacks Project chapter file sites-cohomology.tex; no statement or proof step was rewritten, reordered or removed.\par\medskip
\noindent\textit{Editorial note.} The proof is the author's complete argument; both steps are inside the counted span and the situation cited by the statement is the one introduced earlier in the chapter.\par\medskip
\noindent\textit{Editorial note.} Every labelled reference occurring inside the counted spans is carried in assembly.external\_references as a non-load-bearing navigation pointer (same-chapter ones at their definition site, cross-chapter ones at the invocation inside the counted span), so no load-bearing mathematics is replaced by a stub. No bibliography citation occurs inside the counted spans.\par\medskip
\noindent\textit{Editorial note.} The checked chapter list for same-name collisions is the set of stacks chapters staged in this repository: chow, cohomology, divisors, etale-cohomology, flat, more-algebra, more-morphisms, morphisms, spaces-more-morphisms, varieties together with this card's chapter sites-cohomology.\par\medskip
\noindent\textit{Editorial note.} The preamble is the chapter's own retained minimal preamble (sources/stacks-sites-cohomology/retained-preamble.tex), a verbatim excerpt of the native preamble: the Xy-pic setup and the native theorem declarations with the mathematical macros.\par\medskip
\noindent\textit{Allowed standard prerequisites (editorial).}\begin{itemize}
\item The situation of the Olsson-Laszlo comparison: site, sheaf of rings and weak Serre subcategory with vanishing bounds. Same chapter, the situation introduced earlier in the chapter, which is cited by the counted statement and whose content is the vanishing bound hypothesis used in the proof. Supplies hypotheses (4) and (5) of the statement and the input of the companion results.
\item Bounded version of the equivalence of derived categories of subcategories, and the recovery of objects as derived limits of their truncations. Same chapter, the companion bounded statement and the companion result recovering an object of the subcategory as the derived limit of its truncations. Supplies the two ingredients from which the unbounded statement is assembled.
\end{itemize}\medskip
\ifdefined\hypertarget\hypertarget{sourcepointer-1}{}\fi
\noindent\textit{Source reference.} Original reference 0D6R: sources/stacks-sites-cohomology/pinned-original-source.tex, source lines 5687--5687. Navigation pointer invoked at proof line 6450; the counted span contains the pointer only and the cited result is not reproduced inside it, so the pointer is registered as a non-load-bearing navigation stub in the established form for Stacks material. No bibliography citation is involved.\par\medskip
\ifdefined\hypertarget\hypertarget{sourcepointer-2}{}\fi
\noindent\textit{Source reference.} Original lemma 0D7U: sources/stacks-sites-cohomology/pinned-original-source.tex, source lines 6391--6391. Navigation pointer invoked at proof line 6466; the counted span contains the pointer only and the cited result is not reproduced inside it, so the pointer is registered as a non-load-bearing navigation stub in the established form for Stacks material. No bibliography citation is involved.\par\medskip
\ifdefined\hypertarget\hypertarget{sourcepointer-3}{}\fi
\noindent\textit{Source reference.} Original lemma 0D6U: sources/stacks-sites-cohomology/pinned-original-source.tex, source lines 5796--5796. Navigation pointer invoked at proof line 6467; the counted span contains the pointer only and the cited result is not reproduced inside it, so the pointer is registered as a non-load-bearing navigation stub in the established form for Stacks material. No bibliography citation is involved.\par\medskip
\begin{lemma}
\label{lemma-equivalence-unbounded-one}
\begin{reference}
This is analogous to \cite[Theorem 2.2.3]{six-I}.
\end{reference}
Let $f : (\Sh(\mathcal{C}), \mathcal{O}) \to (\Sh(\mathcal{C}'), \mathcal{O}')$
be a morphism of ringed topoi.
Let $\mathcal{A} \subset \textit{Mod}(\mathcal{O})$
and $\mathcal{A}' \subset \textit{Mod}(\mathcal{O}')$
be weak Serre subcategories. Assume
\begin{enumerate}
\item $f$ is flat,
\item $f^*$ induces an equivalence of categories
$\mathcal{A}' \to \mathcal{A}$,
\item $\mathcal{F}' \to Rf_*f^*\mathcal{F}'$ is an isomorphism
for $\mathcal{F}' \in \Ob(\mathcal{A}')$,
\item $\mathcal{C}, \mathcal{O}, \mathcal{A}$ satisfy the
assumption of Situation \ref{situation-olsson-laszlo},
\item $\mathcal{C}', \mathcal{O}', \mathcal{A}'$ satisfy the
assumption of Situation \ref{situation-olsson-laszlo}.
\end{enumerate}
Then $f^* : D_{\mathcal{A}'}(\mathcal{O}') \to D_\mathcal{A}(\mathcal{O})$
is an equivalence of categories with quasi-inverse given by
$Rf_* : D_\mathcal{A}(\mathcal{O}) \to D_{\mathcal{A}'}(\mathcal{O}')$.
\end{lemma}

\begin{proof}
Since $f^*$ is exact, it is clear that $f^*$ defines a functor
$f^* : D_{\mathcal{A}'}(\mathcal{O}') \to D_\mathcal{A}(\mathcal{O})$
as in the statement of the lemma and that moreover this
functor commutes with the truncation functors $\tau_{\geq -n}$.
We already know that $f^*$ and $Rf_*$ are quasi-inverse
equivalence on the corresponding bounded below categories,
see Lemma \ref{lemma-equivalence-bounded}.
By Lemma \ref{lemma-olsson-laszlo-map-version-one}
with $N = 0$ we see that $Rf_*$ indeed defines a functor
$Rf_* : D_\mathcal{A}(\mathcal{O}) \to D_{\mathcal{A}'}(\mathcal{O}')$
and that moreover this functor commutes with
the truncation functors $\tau_{\geq -n}$.
Thus for $K$ in $D_\mathcal{A}(\mathcal{O})$ the map
$f^*Rf_*K \to K$ is an isomorphism as this is true
on trunctions.
Similarly, for $K'$ in $D_{\mathcal{A}'}(\mathcal{O}')$ the map
$K' \to Rf_*f^*K'$ is an isomorphism as this is true
on trunctions.
This finishes the proof.
\end{proof}

\end{document}
